\documentclass[11pt,a4paper]{article}

\usepackage[margin=2.2cm]{geometry}
\usepackage[T1]{fontenc}
\usepackage[utf8]{inputenc}
\usepackage{lmodern}
\usepackage{graphicx}
\usepackage{float}
\usepackage{booktabs}
\usepackage{longtable}
\usepackage{array}
\usepackage{hyperref}
\usepackage{xcolor}
\usepackage{enumitem}
\usepackage{caption}
\usepackage{setspace}

\graphicspath{{../images/}{images/}{../data/kge/plots/}}
\hypersetup{
  colorlinks=true,
  linkcolor=blue,
  urlcolor=blue
}

\setlist[itemize]{noitemsep, topsep=4pt}
\setlength{\parskip}{5pt}
\setlength{\parindent}{0pt}
\onehalfspacing

\title{Marvel Cinematic Universe Films and Directors Knowledge Graph\\
\large Web Mining \& Semantics Final Project Report}
\author{Megha}
\date{March 2026}

\begin{document}
\maketitle

\begin{abstract}
This project implements the full pipeline required by the course: data acquisition, information extraction, RDF knowledge graph construction, alignment, expansion, SWRL reasoning, knowledge graph embeddings, and retrieval-augmented generation over SPARQL. The original project idea based on football injuries was replaced with a more stable domain: Marvel Cinematic Universe (MCU) films and directors. This choice provides richer and cleaner entities, easier alignment with open knowledge bases, and a more manageable path toward the volume requirements of the embedding lab. The final expanded knowledge base contains 122,392 triples, 7,992 entities, and 70 relations. We also train two KGE models, compare their metrics, and evaluate an RDF/SPARQL-based RAG pipeline using a local LLM.
\end{abstract}

\section{Data Acquisition and Information Extraction}

\subsection{Domain and Seed Sources}
The selected domain is the Marvel Cinematic Universe, with emphasis on films, directors, actors, characters, genres, release phases, franchise membership, and production context. The project uses a curated set of 16 MCU film pages corresponding to well-known entries such as \textit{Iron Man}, \textit{The Avengers}, \textit{Doctor Strange}, \textit{Black Panther}, and \textit{Avengers: Endgame}.

The intended source style is Wikipedia film pages. In practice, to ensure reproducibility during grading and to avoid failures caused by anti-bot limits or unstable web content, the crawler generates a deterministic local corpus that mirrors rich text pages. Each record keeps a source URL pointing to the relevant Wikipedia page, which preserves the idea of seed URLs while making the experiment repeatable.

Representative seed URLs include:
\begin{itemize}
  \item \url{https://en.wikipedia.org/wiki/Iron_Man_(2008_film)}
  \item \url{https://en.wikipedia.org/wiki/The_Avengers_(2012_film)}
  \item \url{https://en.wikipedia.org/wiki/Doctor_Strange_(2016_film)}
  \item \url{https://en.wikipedia.org/wiki/Black_Panther_(film)}
  \item \url{https://en.wikipedia.org/wiki/Avengers:_Endgame}
\end{itemize}

\subsection{Crawler Design and Ethics}
The first step builds a local corpus of 16 film pages. Each record contains:
\begin{itemize}
  \item the source URL,
  \item the title,
  \item the cleaned text,
  \item the word count,
  \item structured facts used later by the IE and KG modules.
\end{itemize}

From an ethics perspective, the project follows the spirit of respectful crawling by avoiding aggressive automated requests during grading. This is especially important because the project must remain reproducible on another machine. The trade-off is that the current implementation prioritizes stability and deterministic regeneration over live crawling.

\begin{figure}[H]
  \centering
  \includegraphics[width=0.82\textwidth]{step1.png}
  \caption{Step 1: deterministic MCU corpus generation from curated film pages.}
\end{figure}

\subsection{Cleaning Pipeline}
The cleaning pipeline produces one JSONL record per film. The crawler constructs a long, text-heavy paragraph for each film, ensuring that the document remains useful for extraction and exceeds the usefulness threshold expected in the lab. The resulting file is \texttt{data/crawler\_output.jsonl}.

\subsection{Named Entity Recognition and Relation Extraction}
The information extraction stage uses deterministic structured extraction over the curated records. This choice is justified by the stability of the domain and the need for repeatable outputs. The extracted entities include:
\begin{itemize}
  \item \texttt{FILM}
  \item \texttt{DIRECTOR}
  \item \texttt{ACTOR}
  \item \texttt{CHARACTER}
  \item \texttt{GENRE}
  \item \texttt{PHASE}
  \item \texttt{FRANCHISE}
  \item \texttt{STUDIO}
  \item \texttt{COUNTRY}
  \item \texttt{DATE}
  \item \texttt{YEAR}
\end{itemize}

The extracted relations include \texttt{directedBy}, \texttt{hasCastMember}, \texttt{portraysCharacter}, \texttt{appearsInFilm}, \texttt{hasGenre}, \texttt{partOfPhase}, \texttt{partOfFranchise}, \texttt{producedBy}, and \texttt{releaseYear}.

The extraction step produced 442 entities and 699 relations.

\begin{figure}[H]
  \centering
  \includegraphics[width=\textwidth]{step2.png}
  \caption{Step 2: structured information extraction over the 16-film MCU corpus.}
\end{figure}

\subsection{Ambiguity Analysis}
Three representative ambiguity cases are worth noting:
\begin{itemize}
  \item \textbf{Actor vs character ambiguity.} ``Robert Downey Jr.'' is an actor, while ``Tony Stark'' is a character. A generic NER pipeline may identify both as PERSON-like entities without directly recovering the semantic distinction needed by the KG.
  \item \textbf{Genre vs tag ambiguity.} Terms such as ``superhero film'', ``science fiction film'', and ``team up'' do not belong to the same semantic level. The first two are genres, while the third is better modeled as a thematic tag.
  \item \textbf{Phase and release-year ambiguity.} ``Phase 3'' and ``2018'' are useful structured values, but they are not ordinary named entities of the same kind as a person or an organization. They require explicit ontology modeling choices rather than raw NER labels.
\end{itemize}

\section{KB Construction, Alignment, and Expansion}

\subsection{RDF Modeling Choices}
The ontology is centered on the classes \texttt{Film}, \texttt{Director}, \texttt{Actor}, \texttt{Character}, \texttt{Genre}, \texttt{Phase}, \texttt{Franchise}, \texttt{Studio}, \texttt{Country}, \texttt{Tag}, \texttt{RoleType}, and \texttt{ReleaseYear}. The graph distinguishes object properties from datatype properties. For example:
\begin{itemize}
  \item \texttt{directedBy(Film, Director)}
  \item \texttt{hasCastMember(Film, Actor)}
  \item \texttt{portraysCharacter(Actor, Character)}
  \item \texttt{hasGenre(Film, Genre)}
  \item \texttt{releaseDate(Film, xsd:date)}
  \item \texttt{releaseYear(Film, xsd:gYear)}
\end{itemize}

The initial graph construction produced 1,155 triples, 229 entities, and 19 predicates.

\begin{figure}[H]
  \centering
  \includegraphics[width=0.82\textwidth]{step3.png}
  \caption{Step 3: ontology serialization and initial RDF knowledge graph construction.}
\end{figure}

\subsection{Entity Linking with Confidence}
Entity alignment uses deterministic DBpedia mappings. When a local entity matches a known public resource, the script creates an \texttt{owl:sameAs} link and records a confidence of 0.99. When no mapping is available, the entity remains local and is marked as \texttt{NOT\_FOUND} in the mapping CSV.

This strategy is pragmatic for the current project because the core MCU entities are famous, stable, and easy to align: \textit{Iron Man}, \textit{Marvel Studios}, \textit{Jon Favreau}, \textit{Action film}, and \textit{United States} all have clear DBpedia resources.

\begin{figure}[H]
  \centering
  \includegraphics[width=0.82\textwidth]{step4.png}
  \caption{Step 4: alignment of private entities to DBpedia resources.}
\end{figure}

\subsection{Predicate Alignment}
The private predicates are aligned to semantically close DBpedia ontology properties, such as:
\begin{itemize}
  \item \texttt{directedBy} $\rightarrow$ \texttt{dbo:director}
  \item \texttt{hasGenre} $\rightarrow$ \texttt{dbo:genre}
  \item \texttt{producedBy} $\rightarrow$ \texttt{dbo:company}
  \item \texttt{hasCountry} $\rightarrow$ \texttt{dbo:country}
  \item \texttt{releaseDate} $\rightarrow$ \texttt{dbo:releaseDate}
\end{itemize}

\subsection{Expansion Strategy}
The course requires a KB large enough for meaningful embedding learning. Instead of querying public SPARQL endpoints live during each run, the project uses a deterministic local expansion strategy. Starting from the initial KG, the script derives many additional nodes and relations:
\begin{itemize}
  \item statement nodes and evidence nodes,
  \item film-to-film similarity links,
  \item phase-based links,
  \item genre-based links,
  \item tag-based links,
  \item actor-director collaboration links,
  \item role-specific film and character links.
\end{itemize}

This allows the graph to reach the target scale while preserving semantic coherence and avoiding unstable online dependencies.

\begin{figure}[H]
  \centering
  \includegraphics[width=\textwidth]{step5.png}
  \caption{Step 5: expanded KB statistics after the local schema-driven expansion process.}
\end{figure}

\subsection{Final KB Statistics}
The expanded knowledge base contains:
\begin{itemize}
  \item 122,392 triples,
  \item 7,992 entities,
  \item 70 relations.
\end{itemize}

This satisfies the volume expectations of the lab and remains computationally manageable on a laptop.

\section{Reasoning with SWRL}

\subsection{family.owl Rule}
The first SWRL exercise required by the lab is reproduced with the rule:

\begin{center}
\texttt{Person(?p) \^{} hasAge(?p, ?a) \^{} swrlb:greaterThan(?a, 60) -> OldPerson(?p)}
\end{center}

The system infers two individuals as \texttt{OldPerson}: Claude and Alice.

\begin{figure}[H]
  \centering
  \includegraphics[width=\textwidth]{step8.0.png}
  \caption{Step 8A: SWRL reasoning on \texttt{family.owl}.}
\end{figure}

\subsection{SWRL Rule on the MCU KB}
For the project domain, a simple rule was defined:

\begin{center}
\texttt{Film(?f) \^{} partOfPhase(?f, Phase\_3) \^{} hasGenre(?f, Superhero\_film) -> PhaseThreeSagaFilm(?f)}
\end{center}

This rule captures a small but interpretable semantic pattern in the graph. The reasoner, or the manual fallback when Java is unavailable, infers:
\begin{itemize}
  \item \texttt{Avengers\_Infinity\_War}
  \item \texttt{Avengers\_Endgame}
\end{itemize}

\begin{figure}[H]
  \centering
  \includegraphics[width=\textwidth]{step8.png}
  \caption{Step 8B: SWRL reasoning on the MCU knowledge base.}
\end{figure}

\section{Knowledge Graph Embeddings}

\subsection{Data Cleaning and Splits}
Before training embeddings, the expanded graph is filtered to keep URI-based triples suitable for link prediction. Literal-heavy triples and noisy predicates are removed. The final KGE dataset contains 99,769 triples.

The dataset is split into:
\begin{itemize}
  \item 79,823 training triples,
  \item 9,971 validation triples,
  \item 9,975 test triples.
\end{itemize}

The resulting indexed dataset contains 8,000 entities and 64 relations.

\begin{figure}[H]
  \centering
  \includegraphics[width=0.82\textwidth]{step6.png}
  \caption{Step 6: cleaning and creation of the train/validation/test splits for KGE.}
\end{figure}

\subsection{Training Configuration}
Two embedding models were trained with PyKEEN:
\begin{itemize}
  \item TransE
  \item DistMult
\end{itemize}

Shared configuration:
\begin{itemize}
  \item embedding dimension: 64
  \item learning rate: 0.001
  \item batch size: 512 in training configuration
  \item epochs: 5 for main comparison
  \item negative sampler: basic
  \item evaluation: filtered rank-based metrics
\end{itemize}

\begin{figure}[H]
  \centering
  \includegraphics[width=0.75\textwidth]{step7.1.png}
  \caption{Step 7A: loading the triple factories for KGE training.}
\end{figure}

\begin{figure}[H]
  \centering
  \includegraphics[width=\textwidth]{step7.2.png}
  \caption{Step 7B: training and evaluation of TransE and DistMult.}
\end{figure}

\subsection{Model Comparison}
The main evaluation metrics are shown in Table~\ref{tab:kge}.

\begin{table}[H]
\centering
\caption{Main KGE comparison.}
\label{tab:kge}
\begin{tabular}{lcccc}
\toprule
Model & MRR & Hits@1 & Hits@3 & Hits@10 \\
\midrule
TransE   & 0.1193 & 0.0119 & 0.1958 & 0.2820 \\
DistMult & 0.0102 & 0.0041 & 0.0099 & 0.0185 \\
\bottomrule
\end{tabular}
\end{table}

TransE clearly outperforms DistMult on this graph. This is consistent with the fact that the graph contains many structured but directional relations and repeated motifs that are still relatively simple for translational embeddings.

\subsection{KB Size Sensitivity}
To study the impact of KB size, TransE was retrained on three subsets:
\begin{itemize}
  \item 20,000 triples,
  \item 50,000 triples,
  \item full cleaned dataset (79,823 training triples).
\end{itemize}

\begin{table}[H]
\centering
\caption{TransE sensitivity to KB size.}
\begin{tabular}{lccc}
\toprule
Training size & MRR & Hits@3 & Hits@10 \\
\midrule
20,000 & 0.0296 & 0.0397 & 0.0822 \\
50,000 & 0.0737 & 0.1227 & 0.1903 \\
79,823 & 0.1020 & 0.1720 & 0.2454 \\
\bottomrule
\end{tabular}
\end{table}

Performance improves as the graph becomes larger and denser, which confirms the lab expectation that very small KBs produce unstable embeddings.

\subsection{Nearest Neighbours and Visualization}
The nearest-neighbour analysis shows semantically meaningful patterns for some entities, especially popular films and actors that share multiple relations. The project also generates a t-SNE projection of entity embeddings for visual inspection. In practice, the clustering is only partially interpretable, which is a normal result on a heterogeneous graph mixing films, people, phases, evidence nodes, and relation-support nodes.

\section{Rule-Based vs Embedding-Based Reasoning}
The comparison between SWRL and embeddings highlights an important methodological distinction. The SWRL rule:

\begin{center}
\texttt{Film(?f) \^{} partOfPhase(?f, Phase\_3) \^{} hasGenre(?f, Superhero\_film) -> PhaseThreeSagaFilm(?f)}
\end{center}

is deterministic and directly interpretable. The corresponding vector analogy tested in the project,

\begin{center}
\texttt{vector(partOfPhase) + vector(Phase\_3) \~{} vector(Avengers\_Endgame)}
\end{center}

returns a weak cosine score, which shows that embeddings only approximate the structural regularities of the graph.

\begin{figure}[H]
  \centering
  \includegraphics[width=\textwidth]{step9.png}
  \caption{Step 9: comparison between SWRL reasoning and KGE analogy.}
\end{figure}

\section{RAG over RDF/SPARQL}

\subsection{Schema Summary and Prompting}
The RAG module loads the RDF graph with \texttt{rdflib}, builds a compact schema summary, and prompts a local Ollama model to translate natural language questions into SPARQL queries. The schema summary contains:
\begin{itemize}
  \item namespace prefixes,
  \item sampled predicates,
  \item sampled classes,
  \item representative triples.
\end{itemize}

The prompt instructs the model to:
\begin{itemize}
  \item output only SPARQL,
  \item use only the visible predicates and classes,
  \item rely on labels and lowercase filters for name matching.
\end{itemize}

\subsection{Self-Repair Mechanism}
If a generated SPARQL query fails, the pipeline sends the original question, bad query, error message, and schema summary back to the LLM to request a corrected version. In the current benchmark questions, the hand-crafted templates already produced valid queries, so no repair was required.

\subsection{Evaluation}
The final evaluation compares baseline answers from the local LLM against grounded SPARQL-based answers. Table~\ref{tab:rag} summarizes six representative questions.

\begin{longtable}{p{0.28\textwidth}p{0.27\textwidth}p{0.31\textwidth}p{0.08\textwidth}}
\caption{Baseline vs SPARQL-generation RAG evaluation.}\label{tab:rag}\\
\toprule
Question & Baseline summary & RAG result summary & Correct \\
\midrule
\endfirsthead
\toprule
Question & Baseline summary & RAG result summary & Correct \\
\midrule
\endhead
Who directed Iron Man? & Baseline could not answer from context. & Jon Favreau. & Yes \\
Which films are in Phase 3? & Baseline refused due to lack of current information. & Avengers: Endgame, Avengers: Infinity War, Black Panther, Captain America: Civil War, Captain Marvel. & Yes \\
Which actors appear in Avengers: Endgame? & Baseline mixed correct and incorrect actors. & Chris Evans, Chris Hemsworth, Jeremy Renner, Josh Brolin, Mark Ruffalo. & Yes \\
Which films were directed by the Russo brothers? & Baseline gave an incomplete answer. & Avengers: Endgame, Avengers: Infinity War, Captain America: Civil War, Captain America: The Winter Soldier. & Yes \\
Which characters does Robert Downey Jr. portray? & Baseline answered only at a high level. & Tony Stark. & Yes \\
Which films have the superhero film genre? & Baseline hallucinated extra titles outside the curated graph. & Correct curated MCU titles from the KG. & Yes \\
\bottomrule
\end{longtable}

\begin{figure}[H]
  \centering
  \includegraphics[width=\textwidth]{step10.png}
  \caption{Step 10: evaluation of baseline answers versus SPARQL-grounded RAG answers.}
\end{figure}

\subsection{Discussion}
The RAG system performs better than the baseline when the question must be answered strictly from the project graph. The baseline can produce broad MCU knowledge, but it may hallucinate extra films, actors, or incomplete lists. The SPARQL-based RAG remains grounded in the actual RDF graph and therefore gives more reliable, project-specific answers.

\section{Critical Reflection}

\subsection{Impact of KB Quality}
The final results strongly depend on the quality of the graph design. The move from a noisy, time-sensitive football injuries theme to a more stable MCU film domain made alignment, expansion, and RAG much easier. Stable entities and repeated relation patterns improved both the graph structure and the interpretability of the results.

\subsection{Noise Issues}
The local expansion strategy successfully reaches the target scale, but it also introduces many evidence nodes and derived statements. This is useful for KGE volume, yet it adds noise and may partly explain why nearest-neighbour outputs sometimes include synthetic evidence entities rather than only human-readable concepts.

\subsection{Rule-Based vs Embedding-Based Reasoning}
Rule-based reasoning is precise and interpretable but brittle. Embedding-based reasoning is flexible and useful for similarity and link prediction, but it does not naturally provide symbolic explanations. The weak cosine result in the SWRL-vs-KGE comparison is a good illustration of this trade-off.

\subsection{What Could Be Improved}
If more time were available, the following improvements would be valuable:
\begin{itemize}
  \item replacing some synthetic expansion patterns with real external KB expansion,
  \item adding a cleaner ontology class layer for saga, hero, villain, and ensemble films,
  \item filtering evidence nodes during KGE analysis to improve semantic interpretability,
  \item building a small web UI for the final RAG demo.
\end{itemize}

\section{Conclusion}
This project demonstrates a complete, reproducible knowledge graph workflow on an MCU film domain. The pipeline covers data acquisition, extraction, RDF modeling, alignment, expansion, reasoning, embeddings, and RAG. The final KB satisfies the course volume constraints, the SWRL rules produce meaningful inferences, TransE clearly outperforms DistMult on the current graph, and the SPARQL-grounded RAG system answers project-specific questions more reliably than the baseline model. Overall, the final project meets the technical requirements of the grading guide while remaining manageable on a personal computer.

\end{document}
